%ALGMEEN
%\documentclass[11pt,a4paper]{article}
%\usepackage{graphicx}
%\usepackage{caption}
%\usepackage{subcaption}
%\usepackage{float}
%\usepackage{ragged2e}
%\usepackage{changepage}
%\usepackage[dutch]{babel}
%\usepackage{subfiles}
%\usepackage[utf8]{inputenc}
%\usepackage{rotating}
%\usepackage{amsmath}  
%\usepackage{tabularx}
%\usepackage{array}
%\usepackage[official]{eurosym}
%\usepackage[toc,page]{appendix}
%\usepackage{nicematrix}
%\usepackage{graphicx}
%\usepackage{tabto}   


%% Style by Bram Gotink, as a clone of the LyX style orginally used, but in clean code.
\documentclass[dutch]{article}
\usepackage[T1]{fontenc}
\usepackage[utf8]{inputenc}
% XeLaTeX/Tectonic negeren inputenc: zet deze tekens om, anders verdwijnen ze
% of worden ze fout gedrukt (§ werd ğ). Onder pdflatex verandert er niets.
\usepackage{iftex}
\ifTUTeX
  \usepackage{newunicodechar}
  \newunicodechar{§}{\S}
  \newunicodechar{‘}{\textquoteleft}
  \newunicodechar{’}{\textquoteright}
  \newunicodechar{“}{\textquotedblleft}
  \newunicodechar{”}{\textquotedblright}
\fi
\usepackage{geometry}
\geometry{verbose,letterpaper,tmargin=2.5cm,bmargin=2.5cm,lmargin=2.5cm,rmargin=2.5cm}
\setlength{\parskip}{\medskipamount}
\setlength{\parindent}{0pt}
\usepackage{textcomp}
\usepackage{titlesec}
\usepackage{enumitem}
\usepackage{color}
\usepackage{ulem}
\usepackage{hyperref}
%\hypersetup{colorlinks=false,urlcolor=cyan,pdfborder=0 0 0}

% VTK Styling

% subsections worden gebruikt als artikels in het HR en Verkiezingsreglement,
% Deze nummering moet doorlopen
\counterwithout{subsection}{section}

% parts zijn "Titel I - De vereniging"
\def\thepart{\Roman{part}}
\def\labelpart{Titel}
\titleformat{\part}
            [hang]
            {\normalfont\huge\bfseries}
            {\labelpart~\thepart~-}
            {10pt}
            {}
% sections zijn "Artikel 1. Naam" enzo
\def\thesection{\arabic{section}}
\def\labelsection{Artikel}
\titleformat{\section}
            [hang]
            {\normalfont\Large\bfseries}
            {\labelsection~\thesection.}
            {10pt}
            {}
% subsections zijn "Artikel 1. Naam", enkel gebruikt in HR
% (daar zijn sections "Afdeling 1 - Naam", zie de .tex van HR)
\def\thesubsection{\arabic{subsection}}
\def\labelsubsection{Artikel}
\titleformat{\subsection}
			[hang]
			{\normalfont\large\bfseries}
			{\labelsubsection~\thesubsection.}
			{10pt}
			{}
% hoogste niveau lijsten zijn §1. etc
\def\theenumi{\S\arabic{enumi}}
\def\labelenumi{\theenumi.}
% tweede niveau wordt eigenlijke eerste niveau
\def\theenumii{\arabic{enumii}}
\def\labelenumii{\theenumii.}
% eerste niveau itemize komt enkel voor in een lijst
\def\labelitemi{}
% layout de lijsten wat mooier
\setlist[enumerate,1]{align=left,labelindent=0pt,labelsep=3em,leftmargin=5em}
\setlist[enumerate,2]{align=left}
\setlist[itemize,1]{align=left,labelindent=0pt,labelsep=3em,leftmargin=5em,labelwidth=0pt}

% locate changes and color them red.
\newenvironment{aanpassing}{\color{red}}{}
\newcommand{\aangepast}[1]{\begin{aanpassing}#1\end{aanpassing}}
\newcommand{\verwijderd}[1]{\aangepast{\sout{#1}}}
% nieuwe (voorgestelde) tekst in het groen
\definecolor{vtkgroen}{rgb}{0,0.55,0}
\newcommand{\toegevoegd}[1]{{\color{vtkgroen}#1}}
\newcommand{\draft}{\textit{\textbf{draft}}}

\usepackage{babel}


\def\labelsection{Afdeling}

\begin{document}



\title{Intern Reglement Vlaamse Technische Kring vzw}
\date{12 mei 2024}

\maketitle


\part{Algemene bepalingen}

    \subsection{Algemene bepalingen}\label{art:algemene-bepalingen}
    \begin{enumerate}
    \item Dit Intern Reglement is geschreven naar aanleiding van de statuten en het Huishoudelijk Reglement van de Vlaamse Technische Kring vzw.
    \item Indien een lid van de Algemene Vergadering dit reglement niet naleeft zonder toestemming van het Bestuursorgaan, kan deze persoon persoonlijk aansprakelijk gesteld worden voor de gevolgen van zijn overtreding.
    \item Verder in dit document wordt de handelsnaam VTK gebruikt voor zowel Vlaamse Technische Kring vzw, als de faculteitskring VTK als het facultair overlegorgaan VTK.
    \end{enumerate}

    % TO HR
    % \subsection{De kringvergadering}
    % \begin{enumerate}
    % \item Het Bestuursorgaan van Vlaamse Technische Kring vzw geldt als de kringvergadering van de faculteitskring VTK, hierna het Praesidium genoemd.
    % \item Bestuurders aangesteld door de Algemene Vergadering die niet verkozen zijn conform de bepalingen van het verkiezingsreglement nemen een adviserende rol waar.
    % \item Het Praesidium is bevoegd voor de vertegenwoordiging van de student op campus (LOKO AV) niveau.
    % \item Dit Intern Reglement vormt eveneens het kringreglement van de faculteitskring VTK, waar nodig aangevuld met het verkiezingsreglement.
    % \end{enumerate}
    
    % \subsection{Het facultair overlegorgaan}
    % \begin{enumerate}
    % \item Het Bestuursorgaan van Vlaamse Technische Kring vzw geldt als de Algemene Vergadering van het facultair overlegorgaan VTK, hierna het FO genoemd.
    % \item Bestuurders aangesteld door de Algemene Vergadering die niet verkozen zijn conform de bepalingen van het verkiezingsreglement nemen een adviserende rol waar.
    % \item Het FO is bevoegd voor het verkiezen van studentenvertegenwoordigers in de opleiding en departementale organen en voor de vertegenwoordiging van de student op facultair (FPOC, faculteitsraad en faculteitsbestuur) en universitair (Stura AV) niveau.
    % \item Dit Intern Reglement vormt eveneens het Intern Reglement van het FO VTK, waar nodig aangevuld met het verkiezingsreglement.
    % \end{enumerate}

\part{Organen en functies}\label{TitelOrganen}

\subsection{Groep 5}

    \subsubsection{Samenstelling}
    \begin{enumerate}
    \item Groep 5 bestaat uit de voorzitter, de ondervoorzitter, de penningmeester(s) en de secretaris van de vzw.
    \end{enumerate}
    
    \subsubsection{Praeses}
    \begin{enumerate}
    \item De rol van Praeses wordt steeds waargenomen door de voorzitter.
    \item De Praeses is de voorzitter van het Praesidium en tevens de voorzitter van het FO.
    \item De Praeses roept de Algemene Vergaderingen, het Bestuursorgaan, de Kringvergadering en de Grocovergadering samen. De praeses zit deze voor en co\"ordineert de werking ervan.
    \item De Praeses stelt in samenspraak met de Grocovergadering de agenda van de Kringvergadering, het Bestuursorgaan en de Algemene Vergaderingen op en is verantwoordelijk voor de uitvoering van de besluiten.
    \end{enumerate}
    
    \subsubsection{Vice-Praeses}
    \begin{enumerate}
    \item De rol van Vice-Praeses wordt steeds waargenomen door de ondervoorzitter.
    \item De Vice-Praeses staat de Praeses bij in al zijn taken en treedt op als zijn vervanger, tenzij de Praeses daartoe een andere persoon aanstelt.
    \item De Vice-Praeses staat ook in voor de interne werking van het Praesidium.
    \item De Vice-Praeses co\"ordineert het medewerkersbeleid van VTK.
    \item De Vice-Praeses is de contactpersoon voor de werkgroepen en jaarwerkingen.
    \item De Vice-Praeses beheert het Salon van VTK. 
    \end{enumerate}
    
    \subsubsection{Financieel Beheerder(s)}
    \begin{enumerate}
    \item De rol van Financieel Beheerder wordt steeds waargenomen door de penningmeester.
    \item De Financieel Beheerder(s) houden nauwgezet de boekhouding van VTK bij.
    \item De Financieel Beheerder(s) zorgen dat de resultaten van de dagelijkse werking worden voorgelegd aan het Bestuursorgaan.
    \item Op de AV's geven de Financieel Beheerder(s) een duidelijk overzicht van de financi\"ele situaties.
    \item De Financieel Beheerder(s) geven de ge\"interesseerde praesidiumleden te allen tijde inzage in de boekhouding.
    \item De Financieel Beheerder(s) zijn verantwoordelijk voor het indienen van de aangiften die verplicht zijn volgens de btw\textendash{}wetgeving.
    \item De Financieel Beheerder(s) zijn verantwoordelijk voor het afwerken van de jaarrekening van hun werkingsjaar.
    \item De Financieel Beheerder(s) stellen minstens tweemaal per werkingsjaar de jaarbegroting voor op het Bestuursorgaan, waarna deze ter goedkeuring wordt voorgelegd op de AV.
    \end{enumerate}
    
    \subsubsection{Secretaris}
    \begin{enumerate}
    \item De rol van Secretaris wordt steeds waargenomen door de secretaris.
    \item De Secretaris staat in voor de ledenadministratie van VTK.
    \item De Secretaris neemt de notulen op van de AV's, het Bestuursorgaan en de Kringvergadering, en geeft deze door aan de leden van de vergadering. Bij afwezigheid van de Secretaris wordt er ad hoc een vervanger aangesteld door de vergadervoorzitter.
    \item De Secretaris co\"ordineert de alumni-werking van VTK %in samenwerking met Alumni Ingenieurs KU Leuven vzw.
    \item De Secretaris beheert het Secretariaat van VTK.
    \item De Secretaris is verantwoordelijk voor het indienen van de aangiften die verplicht zijn volgens de WVV\textendash{}wetgeving.
    \end{enumerate}

\subsection{De Grocovergadering}\label{sec:Groco}

    % \subsubsection{Samenstelling}\label{art:Groco_samenstelling}
    % \begin{enumerate}
    % \item In de Grocovergadering, ook wel Grocomeet genoemd, zetelen de dagelijks bestuurders van VTK.
    % \item Dit omvat de Groepsco\"ordinatoren, hierna Groco's genoemd, en Groep 5.
    % \end{enumerate}
    
    % \subsubsection{Doel en bevoegdheden}\label{art:Groc_doel}
    % \begin{enumerate}
    % \item De Grocovergadering bereidt het Bestuursorgaan en de Kringvergadering voor. De Praeses roept het Bestuursorgaan samen, maar de Grocovergadering stelt de agenda op en is verantwoordelijk voor de uitvoering van de besluiten van het Bestuursorgaan.
    % \item Het dagelijks bestuur wordt waargenomen door de Grocovergadering.
    % \item De Grocovergadering heeft enkel een adviserende rol aan de Kringvergadering met betrekking tot budgettaire beslissingen.
    % \end{enumerate}

    % \subsubsection{Groepsco\"ordinator}\label{art:groco_def}
    \begin{enumerate}
        \item In de Grocovergadering, zoals gespecifieerd in het Huishoudelijk Reglement, zetelen de groepsco\"ordinatoren en Groep 5.
        \item De groepsco\"ordinator, ook wel Groco genoemd, co\"ordineert de werking van de desbetreffende post.
        \item Dit omvat de dagelijkse werking van de post, alsook de vertegenwoordiging van VTK met betrekking tot deze werking.
        \item Daarnaast co\"ordineert de Groco Onderwijs samen met de Praeses de studentenvertegenwoordiging aan de Faculteit Ingenieurswetenschappen en aan de Universiteit en de vertegenwoordiging van VTK met betrekking tot deze studentenvertegenwoordiging, in samenspraak met de betrokken Groco’s. De Groco Onderwijs is ondervoorzitter van het FO.
    \end{enumerate}

\subsection{Het Praesidium}

    \subsubsection{Organisatie}\label{art:praesidium-organisatie}
    \begin{enumerate}
    \item Het Praesidium is onderverdeeld in verschillende posten met verschillende werkomschrijvingen.\label{par:praesidium-posten}
    \item Het Praesidium garandeert te allen tijde minstens één praesidiumlid niet kennelijk onder invloed van alcohol bij ieder onderdeel van de VTK-werking zolang er gewone leden aanwezig kunnen zijn.
    \end{enumerate}

    \subsubsection{Plichten}\label{art:plichten_praesidium}
    \begin{enumerate}
        \item Praesidiumleden zijn verplicht een aantal praesidiumshiften uit te voeren ter ondersteuning van de dagelijkse werking van VTK. Deze shiften worden aan het begin van ieder semester opgesteld en bekendgemaakt door de Vice-Praeses.
        \item Om de continuïteit van de werking te garanderen, wordt ieder praesidiumlid geacht zich te engageren om 2 uur per week te shiften, losstaand van de postwerking.
        \item Er wordt geacht van Praesidiumleden dat ze aanwezig zijn op de Kringvergadering.
    \end{enumerate}
    
    \subsubsection{Praesidiumrekening}
    \begin{enumerate}
    \item De praesidiumrekening is een openstaande rekening die VTK voor zijn praesidiumleden kan voorzien.
    \item De praesidiumrekening wordt beheerd door de Financieel Beheerder(s) van VTK en wordt semestrieel gevorderd. 
    \item Alle praesidiumleden aanvaarden bij hun aanstelling het bestaan en gebruik van deze praesidiumrekeningen.
    \item Het bestaan van praesidiumrekeningen is een bepaalde dienst van VTK naar haar praesidiumleden toe. Deze praesidiumrekeningen kunnen, bij wanbetalers, op elk moment opgeschort worden indien zo beslist door het Bestuursorgaan.
    \item De betalingstermijn voor praesidiumrekeningen bedraagt dertig dagen. Na overschrijden van deze termijn, kan het Bestuursorgaan sancties opleggen.
    \item Afbetalingsplannen of gespreide betalingsmodaliteiten kunnen worden toegekend op eenvoudige aanvraag aan de Financieel Beheerder(s).
    \item Inkom voor activiteiten wordt niet via de praesidiumrekening geregeld.
    \item De stand van de praesidiumrekening kan steeds nagevraagd worden bij de Financieel Beheerder(s).
    \end{enumerate}
    
    \subsubsection{Wedstrijden}\label{art:praesidium-wedstrijden}
    \begin{enumerate}
    \item Leden van het Praesidium kunnen geen wedstrijden georganiseerd door VTK winnen waarin zij door hun lidmaatschap in het voordeel kunnen zijn. Het Bestuursorgaan kan anders beslissen.
    \end{enumerate}
    
    \subsubsection{Beloningen}
    \begin{enumerate}
    \item Het Praesidium wordt niet beloond met medewerkersbonnen voor het post- of praesidiumwerk.
    \item Het Praesidium heeft na afloop van een late of drukke stewardshift recht op een gratis consumptie (bier of frisdrank).
    \item Een praesidiumlid krijgt 3 EUR korting op cantussen (tot een maximum gelijk aan de inkomprijs) indien het praesidiumlid zich extra verantwoordelijk zal gedragen. Dit houdt in dat het praesidiumlid helpt met de volledige afbraak of met stewarden na de cantus.
    Het Bestuursorgaan kan anders beslissen. Voor andere activiteiten kan steeds een voorstel aan het Bestuursorgaan worden voorgelegd.\label{par:CantusKorting}
    \item Een praesidiumlid betaalt voor activiteiten dezelfde inkomprijs als een gewoon lid, eventueel verminderd met een korting zoals in \ref{par:CantusKorting} van dit artikel of een andere korting.
    \item Indien een praesidiumlid een werkshift opneemt tijdens een activiteit (uitgezonderd een tempusshift tijdens een cantus) wordt maar de helft van de inkomprijs na eventuele kortingen zoals beschreven in dit artikel aangerekend aan de kassa. 
    \item Indien een praesidiumlid twee werkshiften opneemt tijdens één activiteit is de inkom gratis voor dit praesidiumlid.
    \item Na afloop van een succesvolle activiteit, hebben alle medewerkers van deze activiteit recht op een gratis consumptie.
    \item Indien een praesidiumlid muziek speelt op een evenement buiten zijn/haar post (DJ in de fak, bandje op galabal) zal dit meetellen als shift op dit evenement of siteshift.
    \end{enumerate}
    
    \subsubsection{Straffen}\label{art:praesidium-straffen}
    \begin{enumerate}
    \item Om het Praesidium bepaalde gedragsregels te laten respecteren kunnen zij gedwongen worden
    hiervoor een straf te ondergaan bepaald door de kringvergadering. Deze straf is uiteraard niet seksueel getint noch vernederend. Een termijn van maximum vier weken tussen het rapporteren aan het Bestuursorgaan en het opleggen van de straf wordt gehanteerd.
    \item Een overtreding van dit Intern Reglement moet gerapporteerd worden aan de Praeses door de getuige(n) of de betrokken personen en dit binnen een termijn van vier weken, met uitzondering van de situatie beschreven in artikel \ref{art:koppel-eerstejaars} waarvoor specifieke maatregelen gelden. In samenspraak met de betrokken personen en de getuige(n) haalt de Praeses dit overtreden op een gepast moment aan op de kringvergadering. Vervolgens beslist de kringvergadering al dan niet om de perso(o)n(en) in kwestie een gepaste straf te geven.
    \item De kringvergadering kan altijd afzien van het recht tot straffen.
    \end{enumerate}
    
    \subsubsection{Praesidiumlint}
    \begin{enumerate}
    \item Een praesidiumlid dat zijn praesidiumlint vrijwillig rondom een andere persoon doet of laat doen of zijn praesidiumlint kwijtgeraakt, wordt voor deze feiten gestraft.
    \end{enumerate}

    \subsubsection{Eerstejaarsstudenten}\label{art:koppel-eerstejaars}
    \begin{enumerate}
    \item De bepaling “een koppel gevormd” slaat op het feit als twee personen elkaar een “Franse kus”
hebben gegeven. 
    Deze bepaling wordt uitgebreid wanneer het koppelvorming tussen een eerstejaarsstudent en een Praesidiumlid betreft. Van zodra er sprake is van enige seksuele handeling is er ook sprake van koppelvorming.
    \item Een praesidiumlid dat een koppel vormt met een eerstejaarsstudent tussen het begin van de onthaaldagen en het einde van het eerste semester, zal gestraft worden.
    \item In geval van koppelvorming met een eerstejaarsstudent wordt de rapporteringstermijn uitgebreid naar het volledige werkingsjaar. 
    \end{enumerate}

\subsection{De Jaarwerkingen}

    \subsubsection{Definitie}
    \begin{enumerate}
    \item Een Jaarwerking is een georganiseerde groep studenten aan de faculteit Ingenieurswetenschappen die zich inzet voor de studenten aan de faculteit Ingenieurswetenschappen van hun studiejaar.
    \item Een Jaarwerking is een onderdeel van VTK, maar beslist meer autonoom over hun werking. 
    \item De Vice-Praeses belegt in het begin van het jaar een vergadering met elke jaarwerking om de rechten en plichten te overlopen en voorziet een algemene uitleg voor de praktische kant hiervan. Op deze vergadering wordt ook een voorzitter van de jaarwerking aangeduid die als contactpersoon fungeert voor de Vice-Praeses.
    \end{enumerate}

    \subsubsection{Plichten}
    \begin{enumerate}
    \item Jaarwerkingen mogen voor activiteiten die door het Bestuursorgaan goedgekeurd zijn, beroep doen op de verzekeringen van VTK vzw. Eventuele extra benodigde verzekeringen worden door VTK vzw afgesloten op vraag van de jaarwerking.
    \item Jaarwerkingen mogen geen sponsoring ontvangen van bedrijven.
    \item De Vice-Praeses krijgt toegang als beheerder op de sociale media van de jaarwerkingen.
    \end{enumerate}

    \subsubsection{Rechten}
    \begin{enumerate}
    \item Jaarwerkingen mogen voor door het Bestuursorgaan goedgekeurde activiteiten beroep doen op de verzekeringen van VTK. Eventuele extra benodigde verzekeringen worden door VTK afgesloten op gegronde vraag van de Jaarwerking.
    \item Jaarwerkingen mogen gebruikmaken van de logistieke diensten van VTK. 
    \item Er staat geen limiet op het gebruik van de voertuigen. De groepscoördinator Logistiek en verantwoordelijke transport hebben echter het recht om onnodig gebruik te wijgeren.
    \item Jaarwerkingen kunnen gebruikmaken van de diensten op Blok 6.
    \end{enumerate}

% To Do: naar HR
% \subsection{De Werkgroepen}

%     \subsubsection{Definitie}
%     \begin{enumerate}
%     \item Een Werkgroep is een georganiseerde groep studenten aan de Faculteit Ingenieurswetenschappen die samen werkt aan iets voor en door de studenten aan de Faculteit Ingenieurswetenschappen.
%     \item Een Werkgroep worden erkend door de Algemene Vergadering van VTK Ondersteuning vzw. 
%     \item De erkende Werkgroepen zijn:
%     \begin{itemize}
%         \item BEST
%         \item Chemix
%         \item Existenz
%         \item IAESTE
%         \item Mechanix
%         \item Statix
%         \item Revue
%     \end{itemize}
%     \item Een Werkgroep is een onderdeel van VTK, maar beslissen volledig autonoom over hun werking. 
%     \item De Vice-Praeses belegt in het begin van het jaar een vergadering met elke werkgroep om de rechten en plichten te overlopen en voorziet een algemene uitleg voor de praktische kant hiervan.
%     \end{enumerate}

%     \subsubsection{Plichten}
%     \begin{enumerate}
%     \item Per Werkgroep zetelen er twee leden in de Algemene Vergadering. Ze worden dan ook aangesteld als werkend lid aan het begin van het werkingsjaar. Deze leden kunnen steeds van het Bestuursorgaan specifieke volmachten krijgen zoals bepaald in de statuten van VTK.
%     \item Werkgroepen moeten een visie op hun beleid ter bekrachtiging voorleggen aan de Algemene Vergadering op de eerste Algemene Vergadering na de overdracht. 
%     \item Werkgroepen kunnen te allen tijde een herziening van deze visie aan het Bestuursorgaan of de Algemene Vergadering voorleggen. 
%     \item VTK controleert of deze visie gevolgd wordt, en kan de werkgroep ter verantwoording roepen 
%     \item Werkgroepen moeten hun volledige begroting ter goedkeuring voorleggen aan de Algemene Vergadering op dezelfde Algemene Vergadering als VTK zijn begroting voorlegt ter goedkeuring.     
%     \item Werkgroepen kunnen te allen tijde nieuwe budget aanvragen aan het Bestuursorgaan of de Algemene Vergadering voorleggen. 
%     \item Werkgroepen moeten een budget per activiteit aan het Bestuursorgaan voorleggen. Wanneer dit punt door het Bestuursorgaan behandeld wordt, wordt een lid van de werkgroep hiervoor uitgenodigd indien zij dit aanvragen of indien de Vice-Praeses en/of het Bestuursorgaan dit nodig achten.
%     \end{enumerate}

%     \subsubsection{Rechten}
%     \begin{enumerate}
%     \item Werkgroepen krijgen een rekeningnummer en bijhorende bankkaarten van VTK tegen 1 september voor het begin van het werkingsjaar. 
%     \item Werkgroepen mogen voor door het Bestuursorgaan goedgekeurde activiteiten beroep doen op de verzekeringen van VTK Ondersteuning vzw. Eventuele extra benodigde verzekeringen worden door VTK Ondersteuning vzw afgesloten op gegronde vraag van de Werkgroep.
%     \item Werkgroepen mogen gebruik maken van de logistieke diensten van VTK. Zij mogen op aanvraag en mits goedkeuring van de Groco Logistiek gebruik maken van de kar tot een limiet van 75 uur, van de auto tot een limiet van 100 uur en het materiaal uit de Kelder.
%     \item Werkgroepen kunnen gebruik maken van de diensten op Blok 6.
%     \item Werkgroepen kunnen gebruik maken van de IT-ondersteuning van VTK, steeds in samenspraak met de Groco IT.
%     \end{enumerate}

\subsection{De Denkgroepen}

    \subsubsection{Definitie}
    \begin{enumerate}
    \item Het Bestuursorgaan of de Algemene Vergadering kan Denkgroepen met een welomlijnde opdracht oprichten om uitgebreidere vraagstellingen te onderzoeken. Ze bepaalt tevens de samenstelling ervan.
    \item Denkgroepen moeten verantwoording afleggen aan het Bestuursorgaan, of indien opgericht door de Algemene Vergadering, aan de Algemene Vergadering. 
    \item Het Bestuursorgaan of, indien van toepassing de Algemene Vergadering, kan bestaande denkgroepen ontbinden.
    \end{enumerate}

% To Do: naar HR
% \subsection{Het Auditcomit\'e}

%     \subsection{Samenstelling}
%     \begin{enumerate}
%     \item Het Auditcomit\'e bestaat uit minstens drie personen, zodat representatie op de Algemene Vergadering verzekerd wordt. Het lidmaatschap van het Auditcomit\'e is onverenigbaar met het lidmaatschap van het Bestuursorgaan en/of het Dagelijks Bestuur.
%     \item De Praeses draagt elk werkingsjaar personen voor voor het Auditcomit\'e. Deze worden door de Algemene Vergadering verkozen voor een termijn van \'e\'en jaar. Zij kunnen onbeperkt herkozen worden.
%     \item De Algemene Vergadering kan steeds een lid van het Auditcomit\'e ontslaan, volgens dezelfde regels als voor de uitsluiting van een Werkend Lid, of een lid van het Auditcomit\'e kan steeds zijn ontslag indienen.
%     \item Het Auditcomit\'e kiest uit hun midden een voorzitter voor het komende jaar. Dit mag niet afdoen aan zijn collegiaal karakter. De voorzitter van het Auditcomit\'e is verantwoordelijk voor het samenroepen van vergaderingen en rapportering aan het Bestuursorgaan. 
%     \item Het Auditcomit\'e kan leden van het Bestuursorgaan uitnodigen voor zijn vergaderingen.
%     \end{enumerate}

%     \subsection{Doel}
%     \begin{enumerate}
%     \item Het Auditcomit\'e heeft tot doel een externe controle te bieden aan het Bestuursorgaan op het gebied van financi\"ele en administratieve verplichtingen. Hierbij handelen ze objectief en zonder de eigenheid van het bestuur te schaden.
%     \item Het Auditcomit\'e is een adviesorgaan. Zijn belangrijkste taak is het informeren van het Bestuursorgaan over zijn verplichtingen en te rapporteren aan de Algemene Vergadering over de uitvoering van deze verplichtingen.
%     \end{enumerate}

%     \subsection{Opdracht}
%     \begin{enumerate}
%     \item Het Auditcomit\'e zal een advies geven over de technische correctheid en haalbaarheid van de begroting. Het Bestuursorgaan bezorgt hiertoe tijdig een exemplaar van de begroting.
%     \item Na de controle op het einde van het eerste semester geeft het Auditcomit\'e opnieuw een advies over de begroting. Hiertoe krijgen zij tijdig de beschikbare resultaten van eerste semester.
%     \item Het Auditcomit\'e zal doorheen het jaar opvolgen of het Bestuursorgaan de boekhouding correct voert en voldaan wordt aan alle verplichtingen (zoals btw, rechtspersonenbelasting, publicaties staatsblad).
%     \item Het Auditcomit\'e zal na afloop van het boekjaar de boekhouding controleren en een advies geven over het al dan niet goedkeuren van de jaarrekening en het al dan niet kwijten van het Bestuursorgaan.
%     \item Het Auditcomit\'e kan steeds in bijzondere gevallen een advies formuleren aan de Algemene Vergadering.
%     \item Elk orgaan van VTK kan ten alle tijde een advies vragen aan het Auditcomit\'e.
%     \item De adviezen vermeld in \S1 en \S4 van dit artikel zijn vereist vooraleer de Algemene Vergadering tot een beslissing kan overgaan, zoals bepaald in Artikel 18 \S1 van de statuten van VTK Ondersteuning vzw. 
%     \item De adviezen/controles vermeld in \S2 en \S4 moeten verplicht aangevraagd/uitgevoerd worden, maar het ontbreken ervan heeft geen gevolgen voor de beslissingsmogelijkheden op de Algemene Vergadering.
%     \end{enumerate}

%     \subsection{Recht van inzage}
%     \begin{enumerate}
%     \item Het Auditcomit\'e heeft het recht van inzage in elk document dat het nodig heeft om zijn opdracht uit te voeren. Zij behandelen deze documenten echter als strikt vertrouwelijk.
%     \end{enumerate}

%     \subsection{Rapportering}
%     \begin{enumerate}
%     \item Het Auditcomit\'e zal steeds een verslag van zijn vergaderingen ter beschikking stellen van de Algemene Vergadering.
%     \end{enumerate}

% \subsection{De Algemene Vergadering}

%     \subsubsection{Samenstelling}\label{art:samenstellingav}
%     \begin{enumerate}
%     \item De Algemene Vergadering is overeenkomstig met de statuten van VTK samengesteld uit de praesidiumleden, twee afgevaardigden van elke erkende Werkgroep en maximaal vijf afgevaardigden van het Auditcomit\'e.
%     \end{enumerate}


% To Do: naar HR
% \part{Draaiboek van de vergaderingen} \label{TitelDraaiboekVergaderingen}

% \subsection{Algemeen Draaiboek}\label{sec:AlgemeenDraaiboek}

%     \subsubsection{Vergadervoorzitter}
%     \begin{enumerate}
%     \item De vergadervoorzitter blijft neutraal gedurende de stemming tenzij bij onbesliste stemming. Indien een stemming onbeslist is, is de stem van de vergadervoorzitter beslissend.
%     \item De vergadervoorzitter mag geen volmacht dragen.
%     \end{enumerate}

%     \subsubsection{Stemming}
%     \begin{enumerate}
%     \item Beslissingen worden genomen bij consensus. Indien er geen consensus bereikt wordt, dient er gestemd te worden. Elk lid van de vergadering kan de stemming aanvragen.
%     \item Iemand kan niet aan een stemming over een agendapunt deelnemen, indien hij/zij niet bij de volledige discussie over dat agendapunt aanwezig was.
%     \item Een stemming gebeurt in principe bij handopsteking, behalve indien anders vermeld in dit reglement. Geheime stemming kan echter ook door \'e\'en van de aanwezige leden aangevraagd worden. Een betrokken natuurlijke persoon woont in geen geval de stemming bij.
%     \item Op iedere vergadering van VTK gebeurt de stemming volgens de principes in dit reglement, tenzij anders vermeld.
%     \item Indien er \'e\'en voorstel voorligt, wordt er gestemd bij absolute meerderheid. Indien er geen absolute meerderheid voor of tegen is, wordt de discussie heropend als minstens \'e\'en aanwezig lid daarom vraagt, en wordt er vervolgens opnieuw gestemd waarbij de onthoudingen niet meetellen.
%     \item Een stemming over \'e\'en voorstel is onbeslist als er bij de tweede stemming over hetzelfde voorstel evenveel voor als tegenstemmen zijn (stemmen voor = stemmen tegen).
%     \item Indien er meerdere voorstellen voorliggen (zowel ordinair als niet ordinair), beslist de vergadervoorzitter over de te volgen stemprocedure.
%     \end{enumerate}

% % To Do: Naar HR
% \subsection{De Grocovergadering}

%     \subsubsection{Samenkomst}
%     \begin{enumerate}
%     \item De Grocovergadering wordt bijeengeroepen door de Praeses. Tevens is de Praeses verplicht de Grocovergadering samen te roepen wanneer minstens \'e\'en vijfde van de leden van de Grocovergadering hierom schriftelijk of elektronisch verzoekt. De datum waarop deze Grocovergadering zal plaatsvinden en de agenda wordt dan in overleg met de betrokken leden door de Praeses bepaald.
%     \end{enumerate}

%     \subsubsection{Besluitvorming}
%     \begin{enumerate}
%     \item Voor stemmingen op de Grocovergadering geldt de neutraliteit van de Praeses niet. 
%     \end{enumerate}
    
%     \subsubsection{Volmachten}
%     \begin{enumerate}
%     \item De leden van de Grocovergadering kunnen zich niet laten vertegenwoordigen door middel van volmachten.
%     \end{enumerate}

%     \subsubsection{Openbaarheid}
%     \begin{enumerate}
%     \item Personen die niet van ambtswege op de Grocovergadering aanwezig moeten zijn, kunnen, in overleg met de vergadervoorzitter, \'e\'en of verscheidene agendapunten van de vergadering bijwonen.
%     \item De notulen worden genomen door een lid van de Grocovergadering en worden ter beschikking gesteld van het Bestuursorgaan.
%     \item Als een Groco zelf niet aanwezig kan zijn op de Grocovergadering kan hij/zij één  iemand van zijn team laten deelnemen in zijn plaats in overleg met de vergadervoorzitter en na instemming van de Grocovergadering.
%     \end{enumerate}

% \subsection{Het Bestuursorgaan}

%     \subsubsection{Samenkomst}
%     \begin{enumerate}
%     \item Het Bestuursorgaan komt in principe elke week op zondag bijeen en telkens wanneer het doel of het belang van VTK het eist. Tevens is de Praeses verplicht het Bestuursorgaan samen te roepen wanneer minstens \'e\'en vijfde van de leden van het Bestuursorgaan hierom schriftelijk of elektronisch verzoekt. De datum waarop deze vergadering zal plaatsvinden en de agenda wordt dan in overleg met de betrokken leden door de Grocovergadering bepaald.
%     \end{enumerate}

%     \subsubsection{Volmachten}
%     \begin{enumerate}
%     \item De leden van het Bestuursorgaan kunnen zich niet laten vertegenwoordigen door middel van volmachten.
%     \end{enumerate}
    
%     \subsubsection{Openbaarheid}
%     \begin{enumerate}
%     \item Personen die niet van ambtswege op het Bestuursorgaan aanwezig moeten zijn, kunnen, in overleg met de vergadervoorzitter en na instemming van het Bestuursorgaan, \'e\'en of verscheidene agendapunten van de vergadering bijwonen.
%     \item Het Bestuursorgaan kan beslissen dat een vergadering of een deel van een vergadering waarop deze personen aanwezig zijn, gesloten wordt verklaard, d.w.z. enkel bijgewoond kan worden door de leden van de vergadering.
%     \end{enumerate}


% To Do: Naar HR
% \part{Studentenvertegenwoordiging}

%     \subsection{Mandaten}
%     \begin{enumerate}
%     \item Bij uitoefening van mandaten als studentenvertegenwoordiger aan de faculteit, universiteit, associatie of daarbuiten is men geen verantwoording verschuldigd aan geen enkel orgaan of persoon binnen VTK indien het vertrouwelijke dossiers betreft, tot het moment dat deze vertrouwelijkheid niet meer vereist is. 
%     \end{enumerate}

% \subsection{Het VTK Bureau}

%     \subsubsection{Samenstelling}
%     \begin{enumerate}
%     \item In het VTK Bureau zetelen alle studenten van de Faculteit Ingenieurswetenschappen, zoals gedefinieerd in Artikel 7 §5 van de statuten van VTK Ondersteuning vzw. 
%     \item De Groco Onderwijs zit het VTK Bureau voor of stelt bij afwezigheid een vervanger aan. Indien er geen vervanger is aangesteld, wordt de vergadering opgeschort.
%     \item Andere personen kunnen, in overleg met de vergadervoorzitter, één of verscheidene agendapunten van de vergadering bijwonen, maar mogen niet deelnemen aan de stemming.
%     \end{enumerate}

%     \subsubsection{Doel en bevoegdheden}
%     \begin{enumerate}
%     \item Het VTK Bureau geldt als de open onderwijsvergadering van VTK. Op deze vergaderingen kunnen onderwerpen uit alle vergaderingen van de departementen, de faculteit, de groep, de campus en de universiteit besproken worden. 
%     \item Het VTK Bureau geeft advies aangaande alle onderwerpen vermeld onder paragraaf 1 van dit artikel.
%     \item Het verslag van het VTK Bureau wordt na goedkeuring op het VTK Bureau, in principe op de eerstkomende vergadering van het Bestuursorgaan voorgelegd ter informatie.
%     \item Alle beslissingen van het VTK Bureau worden aan het Bestuursorgaan ter informatie voorgelegd en besproken onder leiding van de Groco Onderwijs. 
%     \item Over standpunten die niet op het VTK Bureau besproken zijn, beslissen de Praeses en de leden van de post Onderwijs, indien mogelijk aangevuld met de leden van de VTK-delegatie op de desbetreffende vergadering. Deze beslissing wordt toegelicht op het eerstvolgende VTK Bureau.
%     \end{enumerate}

%     \subsubsection{Besluitvorming}\label{art:bureau-besluit}
%     \begin{enumerate}
%     \item In het geval strikte vertrouwelijkheid vereist is, mag het FO bestuur, bestaande uit de leden van de post Onderwijs en de Praeses, een beslissing nemen. Zij zijn hierover geen verantwoording verschuldigd, noch naar het Bestuursorgaan, noch naar het VTK Bureau tot op het moment dat de vertrouwelijkheid niet meer vereist is
%     \end{enumerate}

%     \subsubsection{Openbaarheid}
%     \begin{enumerate}
%     \item De voorzitter van het VTK Bureau kan beslissen dat een vergadering of een deel ervan gesloten wordt verklaard, d.w.z. enkel bijgewoond kan worden door de student-POC'ers, Praeses en de post Onderwijs..
%     \end{enumerate}

% % To Do: naar HR, aanpassen na participatiereglement
% \subsection{Stuver verkiezingen}\label{sec:vertegenwoordiging-verkiezing}

%     \subsubsection{Doel}
%     \begin{enumerate}
%     \item De hieronder beschreven procedure dient ter verkiezing van de studentendelegatie in de verschillende opleidings-POC's, departementsraden en -besturen. 
%     \item Zolang de procedure niet is afgerond blijft de studentendelegatie van het vorige academiejaar actief in zovere zij voldoen aan Artikel 7 §5 van de statuten van VTK Ondersteuning vzw.
%     \end{enumerate}
    
%     \subsubsection{Kandidaatstelling}
%     \begin{enumerate}
%     \item Aan het begin van het academiejaar doet VTK een oproep waarbij studenten aan de faculeit Ingenieurswetenschappen zich kandidaat kunnen stellen voor (een) positie(s) in de verschillende opleidings-POC's, departementsraden en -besturen tot de deadline voor kandidaatstelling.
%     \item Na de deadline voor kandidaatstelling worden de namen van de verschillende kandidaten bekendgemaakt op de website van VTK.
%         \begin{itemize}
%         \item Indien er te weinig of net genoeg kandidaten zijn wordt de procedure zoals in Artikel \ref{art:vertegenwoordiging-stilzwijgend} gevolgd.
%         \item Indien er te veel kandidaten zijn wordt er een compromis gezocht van waarnemers en stemgerechtigden, anders wordt de procedure zoals in Artikel \ref{art:vertegenwoordiging-schriftelijk} gevolgd.
%         \end{itemize}
%     \item De deadline voor kandidaatstelling wordt bepaald door de post Onderwijs.
%     \end{enumerate}
    
%     \subsubsection{Stilzwijgende procedure}\label{art:vertegenwoordiging-stilzwijgend}
%     \begin{enumerate}
%     \item Studenten kunnen bezwaar aantekenen tegen een kandidaat tot de deadline voor bezwaar aantekenen.
%     \item Na de deadline voor bezwaar aantekenen worden
%         \begin{itemize}
%         \item de kandidaten zonder aangetekend bezwaar verkozen als studentenvertegenwoordiger voor het desbetreffende orgaan.
%         \item de kandidaten met aangetekend bezwaar voorgelegd volgens de procedure in Artikel \ref{art:vertegenwoordiging-schriftelijk}.
%         \end{itemize}
%     \item De deadline voor bezwaar aantekenen wordt bepaald door de post Onderwijs.
%     \end{enumerate}
    
%     \subsubsection{Schriftelijke procedure}\label{art:vertegenwoordiging-schriftelijk}
%     \begin{enumerate}
%     \item VTK organiseert een schriftelijke stemming voor de kandidaten waarbij volgende studenten stemgerechtigd zijn:
%         \begin{itemize}
%         \item alle studenten uit de desbetreffende opleiding in het geval van een POC.
%         \item alle studenten uit opleidingen die onder het desbetreffende departement vallen in het geval van een departementsraad of -bestuur. 
%         \end{itemize}
%     \item Kandidaten worden met een gewone meerderheid na afloop van de stemming verkozen.
%     \end{enumerate}

% \subsection{Samenstelling VTK-studentenfractie}

%     \subsubsection{Facultair}
%     \begin{enumerate}
%     \item Voor de samenstelling van de studentenfractie in de faculteitsraad en in de FPOC legt de Groco Onderwijs een lijst ter bekrachtiging voor aan het Bestuursorgaan.
%     \item De samenstelling van de studentendelegaties in de opleidingsPOC’s, departementsraden en departementsbesturen wordt beschreven in Afdeling \ref{sec:vertegenwoordiging-verkiezing}.
%     \item De Praeses en de Groco Onderwijs vertegenwoordigt de studenten in het faculteitsbestuur.
%     \end{enumerate}

%     \subsubsection{Campus}
%     \begin{enumerate}
%     \item  De samenstelling van de stemgerechtigde VTK-delegatie op de LOKO AV wordt beslist door de Praeses.
%     \item De samenstelling van de VTK-delegatie in het geval van een gesloten LOKO AV wordt beslist door de Praeses.
%     \item Bij afwezigheid van de Praeses, worden de bevoegdheden in §1 en §2 uitgeoefend door de Groco Onderwijs. Indien beiden afwezig zijn, wordt een vervanger aangeduid.
%     \end{enumerate}

%     \subsubsection{Universitair}
%     \begin{enumerate}
%     \item  De samenstelling van de stemgerechtigde VTK-delegatie op de Stura AV wordt beslist door de Groco Onderwijs%, rekening houdend met de volgende volgorde van voorrang: De Praeses en de Groco Onderwijs, de leden van de post Onderwijs, de aanwezige praesidiumleden, aanwezige POC’ers, studenten van de faculteit.
%     \item De samenstelling van de VTK-delegatie in het geval van een gesloten Stura AV wordt beslist door de Groco Onderwijs.
%     \item Bij afwezigheid van de Groco Onderwijs, worden de bevoegdheden in §1 en §2 uitgeoefend door de Praeses. Indien beiden afwezig zijn, wordt een vervanger aangeduid.
%     \end{enumerate}


\part{Vastgoed \& rollend materieel}

\subsection{Het Theokot}\label{sec:Theokot}

    \subsubsection{Gebruik en huur}
    \begin{enumerate}
    \item Het gebruik van het Theokot voor VTK-activiteiten georganiseerd door andere organen dan het Theokot zelf is gratis. 
    \item Het gebruik voor praesidium- en ereleden is te allen tijde gratis. \item Voor andere personen en organisaties gelden de huurprijs, waarborg en regels zoals vermeld in het huurcontract dat ter beschikking gesteld wordt door VTK en elk werkingsjaar wordt goedgekeurd door het Bestuursorgaan.
    \item Voor VTK-activiteiten in het Theokot georganiseerd door andere organen dan het Theokot zelf, rekent het Theokot voor drank en snacks de inkoopprijzen aan.
    \item Uitzonderingen op bovenstaande regels worden door het Bestuursorgaan goedgekeurd.
    \item In het Theokot geldt een algemeen rook\textendash{} en vapeverbod.
    \end{enumerate}

% To Do: nota huur toevoegen
\subsection{Fakbar 't ElixIr}\label{sec:Elixir}

    \subsubsection{Gebruikt en huur}
    \begin{enumerate}
        \item Het gebruik van Fakbar 't ElixIr door andere posten van het VTK-praesidium is gratis. Winst of verlies gemaakt op deze activiteiten is voor Fakbar 't ElixIr.
        \item Indien er een activiteit van een post van VTK in samenwerking met een andere kring wordt georganiseerd, bedraagt de huur 100 EUR incl. btw. Winst of verlies van deze activiteit is ten laste van de desbetreffende post en kring. 
        \item Werkgroepen en jaarwerkingen van VTK hebben het recht om eenmaal per semester, mits goedkeuring van de groepsco\"ordinator Fakbar, Fakbar 't ElixIr te huren. Hiervoor wordt een huur van 100 EUR incl. btw aangerekend. De winst of verlies die op deze avond gemaakt wordt is voor de desbetreffende werkgroep/jaarwerking.
        \item Verhuur aan derden is geregeld zoals vermeld in het huurcontract, dat ter beschikking gesteld wordt door VTK en elk werkingsjaar door de kringvergadering moet worden goedgekeurd.
        \item Verhuur aan (oud\textendash{})postleden van de post Fakbar is te allen tijde gratis. Deze verhuur kan enkel gebeuren voor een priv\'e-evenement op momenten buiten de reguliere openingsuren.
        \item Uitzonderingen op bovenstaande regels worden door de kringvergadering goedgekeurd.
        \item In Fakbar 't ElixIr geldt een algemeen rook\textendash{} en vapeverbod.
    \end{enumerate}


\subsection{Camerabewaking Fakbar 't ElixIr}
    \begin{enumerate}
        \item In Fakbar 't ElixIr zijn camera's ge\"installeerd ter bewaking van de algemene orde en veiligheid. 
        \item Deze camera's worden enkel gebruikt voor de doeleinden hier beschreven.
    \end{enumerate}
    \subsubsection{Opslag camerabeelden}
    \begin{enumerate}
        \item Beelden gemaakt door de camerabewaking blijven bewaard voor 14 dagen, en zullen na deze termijn automatisch verwijderd worden, tenzij ze het voorwerp zijn van vermeende feiten.
        \item Camerabeelden worden nooit handmatig verwijderd voor deze termijn verstreken is.
    \end{enumerate}

    \subsubsection{Algemene provisies}
    \begin{enumerate}
        \item De camerabeelden mogen enkel bekeken worden door de personen die hiervoor in dit reglement gemachtigd zijn. Gemachtigden hebben niet het recht andere mensen toegang te verlenen tot de beelden. Uitzonderingen op de gemachtigden zullen steeds door de Kringvergadering moeten worden goedgekeurd.
        \item Gemachtigden hebben de verantwoordelijkheid om ervoor te zorgen dat het nooit mogelijk is voor derden om mee te kijken op het moment dat de beelden bekeken worden. 
        \item Er wordt altijd duidelijk gecommuniceerd aan alle aanwezigen in 't ElixIr wie de avondverantwoordelijke is. Per avond is dit altijd maar \'e\'en persoon.
        \item De inhoud van de beelden mag alleen vrij besproken worden tussen gemachtigden, zonder dat dit gemachtigden automatisch gemachtigd maakt. Wanneer gerapporteerd moet worden aan de Kringvergadering mag de inhoud van de beelden aan de Kringvergadering vrijgegeven worden.
    \end{enumerate}

    \subsubsection{Toegang tot de beelden}
    \begin{enumerate}
        \item De camerabeelden worden geraadpleegd door de volgende personen slechts op de daarvoor toegestane tijdstippen en op het moment dat ze gemachtigd zijn:
        \begin{itemize}
            \item Leden van het Fakbarteam
            \item Groepsco\"ordinator Fakbar
            \item Praeses
        \end{itemize}
        \item Enkel wanneer gemachtigd hebben de gemachtigden toegang tot de camerabeelden die betrekking hebben op het voorval.
        \item Betrokken personen kunnen toegang krijgen tot stilstaande beelden, indien er hiervoor goedkeuring is gegeven door een gemachtigde.
    \end{enumerate}

    \subsubsection{Raadplegen van de beelden}
    \begin{enumerate}
        \item Enkel indien de toegang tot de camerabeelden tot een onmiddellijke verbetering van de situatie kan leiden is de avondverantwoordelijke gemachtigd om beelden te bekijken op de avond zelf, voor de beelden die gerelateerd zijn aan, maar niet beperkt tot de volgende feiten:
        \begin{itemize}
            \item Diefstal
            \item Seksueel grensoverschrijdend gedrag
            \item Vandalisme
            \item Geweld
        \end{itemize}
        De avondverantwoordelijke meldt dit achteraf aan de groepsco\"ordinator.
        \item Voor de vermelde feiten is de avondverantwoordelijke op een later tijdstip gemachtigd om de camerabeelden te bekijken, mits hij hiervoor de groepsco\"ordinator op voorhand inlicht. 
        \item De beelden worden op de avond zelf zo weinig mogelijk bekeken. Voor diefstal wordt er enkel naar de beelden gekeken indien de afwezigheid van het gestolen voorwerp de lichamelijke integriteit van de persoon zou kunnen compromitteren.
        \item\label{art:groepscoo_camera_fak} De groepsco\"ordinator bekijkt de beelden, buiten in hoedanigheid als avondverantwoordelijke, enkel in de volgende gevallen:
            \begin{itemize}
                \item Klachten over de avondverantwoordelijke van een avond
                \item Vermoeden van eerdergenoemde feiten waar een praesidiumlid bij betrokken is
                \item Zaken die vallen onder de meldplicht van de avondverantwoordelijke
            \end{itemize}
        \item Wanneer de groepsco\"ordinator zelf avondverantwoordelijke was, gaat de verantwoordelijkheid over de controle van de avondverantwoordelijke en meldplicht naar de praeses.
        \item De praeses bekijkt de beelden enkel in de volgende gevallen:
        \begin{itemize}
            \item Klachten over de groepsco\"ordinator
            \item Zaken die vallen onder de meldplicht van de groepsco\"ordinator.
        \end{itemize}
    \end{enumerate}

    \subsubsection{Rapportering en meldplicht}
    \begin{enumerate}
        \item Wanneer camerabeelden bekeken worden door een gemachtigde, heeft deze schriftelijke meldplicht van het bekijken van deze beelden bij de andere gemachtigden. Deze gegevens worden bijgehouden.
        \item Als de praeses de camerabeelden bekijkt, meldt hij dit aan de Kringvergadering.
        \item Als er blijkt dat een klacht als (gedeeltelijk) gegrond wordt bevonden door de avondverantwoordelijke bespreekt de groepsco\"ordinator dit met de praeses. 
        \item Voor klachten over de praeses ligt de verantwoordelijkheid bij de groepscoördinator. Deze rapporteert hierover rechtstreeks aan de kringvergadering.
        \item De camerabeelden worden nooit verstuurd naar derden met uitzondering van de politie en/of onder gerechtelijke bevel.
        \item Bij het overdragen van camerabeelden aan de politie wordt de kringvergadering daarvan altijd op de hoogte gesteld op de eerstvolgende vergadering.
        \item Alle gemachtigden hebben controleplicht over de andere gemachtigden. De gemachtigden, exclusief de groepscoördinator en de praeses, moeten rapporteren aan de groepscoördinator en de praeses, tenzij deze beiden het voorwerp van een inbreuk zijn. Indien deze wel het voorwerp zijn van een inbreuk moet de gemachtigde rechtstreeks aan de kringvergadering rapporteren. 
        \item De groepsco\"ordinator en de praeses moeten over inbreuken van andere gemachtigden rapporteren aan de kringvergadering.
    \end{enumerate}
    
    \subsubsection{Inbreuken}
    \begin{enumerate}
        \item Elk praesidiumlid heeft meldplicht bij de groepscoördinator bij het opmerken van een inbreuk op dit reglement. Indien de inbreuk gaat over groepscoördinator heeft het praesidiumlid meldplicht bij de praeses. Indien de inbreuk gaat over de praeses heeft het praesidiumlid meldplicht bij de groepscoördinator. Indien beiden het voorwerp zijn meldt het praesidium dit op de kringvergadering.
        \item Bij een inbreuk op dit reglement waar de praeses voorwerp van is, wordt dit op de eerstvolgende kringvergadering behandeld door persoon op de kringvergadering met de hoogste anciënniteit. Indien er meerdere personen zijn met dezelfde anciënniteit zijn wordt de oudste persoon van de groep mensen met de hoogste anciënniteit aangesteld.  
    \end{enumerate}

    \subsubsection{Alarmbelprocedure}
    \begin{enumerate}
        \item De kringvergadering heeft altijd het recht om een of meerdere vertegenwoordiger(s) met een specifiek mandaat in haar naam aan te stellen die gemachtigd worden om de camerabeelden te bekijken binnen het mandaat. Deze vertegenwoordiger(s) heeft/hebben enkel rapporteringsplicht aan de kringvergadering. 
    \end{enumerate}

\subsection{Blok 6}

    \subsubsection{Het Secretariaat}\label{art:secri}
    \begin{enumerate}
    \item VTK stelt in haar Secretariaat een telefoontoestel, een kopietoestel en computers ter beschikking. Er wordt altijd voorrang verleend aan VTK-werk bij het gebruik van deze apparaten.
    \item Het Secretariaat en het materiaal aanwezig wordt beheerd door de Secretaris. De computers worden beheerd door IT.
    \item Roken, vapen, het laten rondslingeren van eten en het ongeoorloofd dumpen van materiaal in het Secretariaat is ten strengste verboden.
    \end{enumerate}

    \subsubsection{Het Salon}\label{art:salon}
    \begin{enumerate}
    \item VTK stelt in haar Salon een vergaderruimte ter beschikking. Het Salon kan gebruikt worden voor vergaderingen i.v.m. VTK, indien vooraf aangevraagd aan de Vice-Praeses.
    \item Het Salon wordt na elk gebruik netjes achtergelaten.
    \item Het Salon en het materiaal aanwezig wordt beheerd door de Vice-Praeses.
    \item Roken, vapen, het laten rondslingeren van eten en het ongeoorloofd dumpen van materiaal in het Salon is ten strengste verboden.
    \end{enumerate}

    \subsubsection{Opslagplaatsen}\label{art:kelder}
    \begin{enumerate}
    \item VTK stelt een of meerdere opslagruimte(s) ter beschikking voor VTK-doeleinden.
    \item De opslagruimte(s) en het materiaal aanwezig wordt beheerd door de post Logistiek.
    \item Roken, vapen, het laten rondslingeren van eten en het ongeoorloofd dumpen van materiaal in de opslagruimte(s) is ten strengste verboden.
    \end{enumerate}

    \subsubsection{Gemeenschappelijke ruimtes Blok 6/0\&1}\label{art:blok6}
    \begin{enumerate}
    \item De gangen, sanitaire ruimten, keukens, de inkomhal, de omgeving en al de overige mogelijke bergplaatsen van Blok 6/0\&1 worden niet gebruikt om materiaal te stockeren.
    \end{enumerate}


\subsection{Rollend materieel}

    \subsubsection{De Kar en Auto}
    \begin{enumerate}
    \item VTK stelt met de Kar en de Auto voertuigen ter beschikking om de werking te ondersteunen.
    \end{enumerate}

    \subsubsection{De chauffeurs}
    \begin{enumerate}
    \item De kringvergadering duidt elk werkingsjaar een aantal permanente karchauffeurs aan. Enkel zij hebben het recht om de Kar te besturen, tenzij expliciete goedkeuring van een uitzondering op deze regel door de Groco Logistiek of de kringvergadering wordt gegeven. 
    \item De kringvergadering duidt elk jaar een aantal permanente autochauffeurs aan. Enkel zij hebben het recht om de Auto te besturen, tenzij expliciete goedkeuring van een uitzondering op deze regel door de Groco Logistiek of de kringvergadering wordt gegeven.
    \item De karchauffeurs mogen automatisch ook de Auto besturen.
    \item De Groco Logistiek bepaalt in het begin van het jaar hoeveel autochauffeurs elke andere post kan voordragen gebaseerd op het gebruik van de auto door die post het jaar voordien. Ook elke erkende werkgroep en jaarwerking van VTK kan minstens \'e\'en autochauffeur voordragen.
    \item Chauffeurs dienen te allen tijde te beschikken over een geldig rijbewijs B, en dit sedert \'e\'en jaar. Uitzonderingen hierop dienen steeds goedgekeurd te worden door de kringvergadering.
    \item De kringvergadering kan de chauffeurs steeds wijzigen.
    \item De eerste boete die een chauffeur oploopt tijdens het rijden voor VTK, kan voor 50\% worden terugbetaald aan de chauffeur. Indien de financieel beheerders niet akkoord gaan met de terugbetaling, kunnen ze de kwestie aanhalen op de eerstvolgende kringvergadering, waar beslist wordt om de boete al dan niet terug te betalen. 
    \end{enumerate}

    \subsubsection{Bijrijders}
    \begin{enumerate}
    \item De chauffeur moet bij elke rit met de Kar of Auto verplicht vergezeld worden van een bijrijder.
    \end{enumerate}

    \subsubsection{Gebruik van de Kar en Auto}
    \begin{enumerate}
    \item De Kar en Auto worden enkel gebruikt voor VTK-doeleinden. Eventuele
    uitzonderingen op deze regel kunnen toegestaan worden door de Groco Logistiek. VTK-doeleinden hebben steeds voorrang.
    \item Het tarief voor het gebruik voor niet-VTK-doeleinden is 7,50 EUR incl. btw per uur. 
    \item Het tarief voor het gebruik door praesidiumleden is 5,00 EUR incl. btw per uur. De kringvergadering kan hier anders over beslissen.
    \item Bij korte onbeschikbaarheid van de Kar of Auto, kan er met priv\'evoertuigen voor VTK-doeleinden gereden worden. Indien vooraf aangevraagd aan de Groco Logistiek, vergoedt VTK de priv\'ekilometers aan de tarieven die elk werkingsjaar door de kringvergadering worden bepaald.
    \item Bij lange onbeschikbaarheid van de Kar of Auto zorgt de kringvergadering voor een vervangvoertuig.
    \end{enumerate}

\subsection{Andere Voorzieningen}

    \subsubsection{Materiaal}
    \begin{enumerate}
    \item Logistiek materiaal kan na voorafgaandelijke aanvraag aan
    de post Logistiek ontleend worden. Opnieuw heeft het gebruik voor VTK-doeleinden altijd voorrang. De post Logistiek kan op elk moment ontleend materiaal terugvorderen.
    \item Het ontlenen en huren van logistiek materiaal door leden, andere personen en organisaties is onderhevig aan de voorwaarden zoals vermeld in het huurcontract dat ter beschikking gesteld wordt door VTK.
    \item De gevraagde huurprijzen en waarborgen worden in het begin van het werkingsjaar vastgelegd door de kringvergadering.
    \item Alle materiaal dat VTK bezit en buiten de verantwoordelijkheid van de Groco Logistiek valt kan uitgeleend worden mits toestemming van de kringvergadering of van de Groco onder wiens verantwoordelijkheid het materiaal valt.
    \end{enumerate}


\part{Medewerkers}

    \subsection{Definitie shiften}
    \begin{enumerate}
    \item Een shift is een of meerdere taken die uitgevoerd worden door een persoon onder verantwoordelijkheid van een praesidiumlid.
    \item Voor de medewerkersbeloningen tellen enkel shiften op activiteiten georganiseerd door VTK en niet door jaarwerkingen of werkgroepen.
    \item Op grote evenementen of evenementen van lange duur voorziet de Vice-Praeses medewerkerseten en medewerkersdrank voor alle aanwezige shifters inclusief de praesidiumshifters.
    \end{enumerate}

    \subsection{Gewone en vaste medewerkers}
    \begin{enumerate}
    \item Elk gewoon lid dat minstens één shift heeft gedaan in het huidige werkingsjaar geldt als medewerker van VTK. Leden van het Bestuursorgaan kunnen geen medewerker zijn.
    \item Een gewone medewerker is een medewerker die geen vaste medewerker is.
    \item Vaste medewerkers zijn medewerkers die ten minste 15 shiften gewerkt hebben gedurende een werkingsjaar.
    \item Het aantal gewerkte shiften kan niet overgedragen worden naar een volgend werkingsjaar. Aan het begin van elk werkingsjaar start iedereen dus met nul shiften. Voor lange shiften (meer dan 2 uur) geldt dat elk voltooid blok van twee uur, meetelt als één shift voor het statuut van vaste medewerker.
    \item Vaste medewerkers hebben recht op volgende voordelen:
        \begin{enumerate}
        \item inkomkaarten kunnen op voorhand bij de Vice-Praeses besteld worden,
        \item een klein VTK-aandenken, en
        \item toegang tot het medewerkersfeestje op het einde van het jaar.
        \end{enumerate}
    \item Het Bestuursorgaan kan steeds mensen die zich verdienstelijk hebben ingezet voor VTK tot vaste medewerkers benoemen.
    \item Een medewerker krijgt een bedrag, gelijk aan de helft van de betaalde inkomprijs, terugbetaald per werkshift tijdens cantussen (geen tempus- of plaspolitieshiften) en TD’s, met een maximum van de betaalde inkomprijs. 
    \end{enumerate}

    \subsection{Medewerkersbonnen}
    \begin{enumerate}
    \item Per shift voor 02u00 krijgt een medewerker één medewerksbon per uur. Per werkshift na 02u00 krijgt een medewerker 1.5 medewerkersbonnen per uur (afgerond door de Vice-Praeses).
    \item Medewerkersbonnen kunnen worden afgehaald op het kot van de Vice-Praeses. De Vice-Praeses geeft duidelijk aan wanneer hij/zij aanwezig is op kot.
    \item In 't ElixIr kan met 1 bonnetje een standaardconsumptie besteld worden. Met 2 bonnetjes kan men in het Theokot een standaardbroodje bestellen.
    \item Standaardconsumpties en -broodjes worden bepaald door het Bestuursorgaan.
    \item Er wordt geen geld teruggeven op een betaling met medewerkersbonnen.
    \end{enumerate}


\part{Aanvullingen}

    \subsection{Ereleden}
    \begin{enumerate}
    \item De ereleden verkrijgen op de VTK-activiteiten een korting van 5 EUR met een maximum van de inkomprijs op vertoon van hun erelidkaart, zoals vastgelegd door het Bestuursorgaan. Indien noodzakelijk moet de korting vooraf aangevraagd worden.
    \item Het galabal en cantussen zijn gratis toegankelijk voor alle ereleden.
    \item Uitzonderingen op deze regels worden aan het Bestuursorgaan voorgelegd.
    \item Voor cantussen kan een maximumpercentage ereleden worden opgelegd door het Bestuursorgaan teneinde de gewone leden niet te benadelen.
    \item Elk praesidiumlid krijgt per jaar praesidium een erelidmaatschap van 1 academiejaar gratis. Het praesidiumlid mag zelf beslissen wanneer hij de jaren erelidmaatschap opneemt, zolang deze binnen de 5 jaar na het afstuderen opgenomen worden en het praesidiumlid zijn/haar praesidiumrekening betaald heeft. Het Bestuursorgaan kan anders beslissen.
    \end{enumerate}


\part{Slotbepalingen}

    \subsection{Wijzigingen}
    \begin{enumerate}
    \item De bepalingen van dit Intern Reglement worden goedgekeurd en eventueel gewijzigd door een gewone meerderheid van het Bestuursorgaan van de Vlaamse Technische Kring vzw.
    \end{enumerate}

    \subsection{Historiek der wijzigingen}
    \begin{enumerate}
    \item Dit Huishoudelijk Reglement werd aangenomen op de Algemene Vergaderingen van VTK vzw en VTK Ondersteuning vzw op 3 januari 2007.
    \item Dit Huishoudelijk Reglement werd gewijzigd op de Algemene Vergaderingen van VTK vzw en VTK Ondersteuning vzw op 9 december 2007.
    \item Dit Huishoudelijk Reglement werd gewijzigd op de Algemene Vergaderingen van VTK vzw en VTK Ondersteuning vzw op 16 november 2008. 
    \item Dit Huishoudelijk Reglement werd gewijzigd op de Algemene Vergaderingen van VTK vzw en VTK Ondersteuning vzw op 29 november 2009.
    \item Dit Huishoudelijk Reglement werd gewijzigd op de Algemene Vergaderingen van VTK vzw en VTK Ondersteuning vzw op 28 november 2010.
    \item Dit Huishoudelijk Reglement werd gewijzigd op de Algemene Vergaderingen van VTK vzw en VTK Ondersteuning vzw op 18 november 2012.
    \item Dit Huishoudelijk Reglement werd gewijzigd op de Algemene Vergaderingen van VTK vzw en VTK Ondersteuning vzw op 19 oktober 2014.
    \item Dit Huishoudelijk Reglement werd gewijzigd op de Algemene Vergaderingen van VTK vzw en VTK Ondersteuning vzw op 8 november 2015.
    \item Dit Huishoudelijk Reglement werd gewijzigd op de Algemene Vergadering van VTK Ondersteuning vzw op 1 mei 2016 naar aanleiding van de vereffening van VTK vzw.
    \item Dit Huishoudelijk Reglement werd gewijzigd op de Algemene Vergadering van VTK Ondersteuning vzw op 25 februari 2018.
    \item Dit Huishoudelijk Reglement werd gewijzigd op de Algemene Vergadering van VTK Ondersteuning vzw op 13 oktober 2019.
    \item Dit Huishoudelijk Reglement werd gewijzigd op de Algemene Vergadering van VTK Ondersteuning vzw op 13 juli 2020.
    \item Dit Intern Reglement werd gewijzigd op de Algemene Vergadering van VTK Ondersteuning vzw op 19 mei 2021.
    \item Dit Intern Reglement werd gewijzigd op het Bestuursorgaan van VTK Ondersteuning vzw op 1 juli 2021.
    \item Dit Intern Reglement werd gewijzigd op het Bestuursorgaan van Vlaamse Technische Kring vzw op 12 mei 2024.
    \end{enumerate}

\end{document}